\documentclass[12pt,a4paper]{article}
\usepackage[utf8]{inputenc}
\usepackage{geometry}
\geometry{a4paper, margin=1in}
\usepackage{graphicx}
\usepackage{hyperref}
\usepackage{listings}
\usepackage{xcolor}
\usepackage{amsmath}
\usepackage{tabularx}
\usepackage{booktabs}
\usepackage{float}
\usepackage{tcolorbox}
\usepackage{parskip}

\definecolor{codegreen}{rgb}{0,0.6,0}
\definecolor{codegray}{rgb}{0.5,0.5,0.5}
\definecolor{codepurple}{rgb}{0.58,0,0.82}
\definecolor{backcolour}{rgb}{0.95,0.95,0.92}

\lstdefinestyle{mystyle}{
    backgroundcolor=\color{backcolour},   
    commentstyle=\color{codegreen},
    keywordstyle=\color{blue},
    numberstyle=\tiny\color{codegray},
    stringstyle=\color{codepurple},
    basicstyle=\ttfamily\footnotesize,
    breakatwhitespace=false,         
    breaklines=true,                 
    captionpos=b,                    
    keepspaces=true,                 
    numbers=left,                    
    numbersep=5pt,                  
    showspaces=false,                
    showstringspaces=false,
    showtabs=false,                  
    tabsize=2
}
\lstset{style=mystyle}

\title{\Huge \textbf{Comprehensive System Architecture \& Code Explanation} \\ \vspace{0.5cm} \Large Smart Air Guardian Project}
\author{Engineering \& Development Team}
\date{\today}

\begin{document}

\maketitle
\tableofcontents
\newpage

\section{Executive Summary}
The \textbf{Smart Air Guardian} is a state-of-the-art IoT air quality monitoring and medical-grade diagnostic system. It bridges the gap between low-level hardware sensors, edge computing, cloud telemetry, and progressive web application interfaces. This document serves as the definitive, professional-grade explanation of the entire system architecture, the Machine Learning pipeline, and a granular breakdown of the C++ firmware running on the ESP32 edge node.

\section{System Architecture Overview}
The system is divided into three primary technological pillars:
\begin{enumerate}
    \item \textbf{The Hardware Edge Node (ESP32)}: Acts as the physical data acquisition and edge inference engine.
    \item \textbf{The Cloud Broker (Firebase RTDB)}: Acts as the ultra-low latency middleware connecting the hardware to the global internet.
    \item \textbf{The Command Center (PWA)}: The user-facing dashboard built with HTML, CSS, and JS, featuring a premium Claymorphic design.
\end{enumerate}

\section{Hardware and Sensor Matrix}
The ESP32 microcontroller continuously polls the following environmental sensors:
\begin{itemize}
    \item \textbf{DHT22}: Measures ambient Temperature and Humidity (Digital GPIO 4).
    \item \textbf{Sharp GP2Y1010}: An optical dust sensor measuring PM2.5 particulate matter (Analog ADC1\_CH4 - Pin 32).
    \item \textbf{MQ-135}: Gas sensor detecting general VOCs and CO2 equivalents (Analog ADC1\_CH6 - Pin 34).
    \item \textbf{MQ-7}: Highly sensitive Carbon Monoxide (CO) sensor (Analog ADC1\_CH7 - Pin 35).
    \item \textbf{Actuators}: Active Buzzer (Pin 26) and RGB LED Channels (Pins 25, 27, 12).
\end{itemize}

\section{Machine Learning Pipeline and Edge AI}
The ESP32 runs Edge AI, meaning it makes complex predictions without requiring an internet connection. This is achieved by training classical ML models in Python and exporting the final mathematical formulas to C++.

\subsection{Data Acquisition and Preparation}
Python scripts (\texttt{download\_all\_data.py}) download historical records of PM2.5, CO, Temperature, and Humidity in Bangladesh, saving them as \texttt{Master\_AQI\_Weather.csv}.

\subsection{Training the AI Models (\texttt{train\_final\_models.py})}
The \texttt{scikit-learn} library is utilized to train three distinct models based on the master dataset:

\subsubsection{1. Linear Regression (Forecaster)}
\textbf{Goal:} Predict the exact quantitative AQI 1 hour into the future.\\
\textbf{Mechanism:} The algorithm draws a line of best fit through the 4-dimensional historical data. It determines specific weight coefficients for PM2.5, CO, Temperature, and Humidity. 
\\
\textbf{C++ Implementation:} The ESP32 multiplies the real-time scaled inputs by these weights and adds an intercept to yield the forecasted AQI.

\subsubsection{2. Logistic Regression (Hazard Sentinel)}
\textbf{Goal:} Output a binary probability (0\% to 100\%) indicating an imminent emergency hazard (AQI $>$ 150).\\
\textbf{Mechanism:} Logistic regression calculates a mathematical $Z$ score, and then passes it through a \textbf{Sigmoid Activation Function} to restrict the output strictly between $0.0$ and $1.0$.
\\
\textbf{C++ Implementation:} 
\begin{lstlisting}[language=C++]
float z = (23.5640f * s_pm25) + (-0.3901f * s_co) + (-0.0083f * s_temp) + (-0.0110f * s_humid) - 4.4767f;
float prob = 1.0f / (1.0f + exp(-z)); // Sigmoid activation
\end{lstlisting}

\subsubsection{3. Decision Tree Classifier (Risk Assessor)}
\textbf{Goal:} Categorize air into 4 clinical risk buckets: Good (0), Moderate (1), Unhealthy (2), Hazardous (3).\\
\textbf{Mechanism:} The algorithm constructs a flowchart based on optimal threshold splits (Gini impurity).\\
\textbf{C++ Implementation:} Translated into a highly efficient series of nested \texttt{if/else} statements in the firmware, executing in nanoseconds.

\subsection{Feature Normalization (Scaling)}
For the ML formulas to work, numbers must be constrained. The Python script calculates a \texttt{MinMaxScaler} to find the absolute maximum and minimum historical values. The ESP32 utilizes the \texttt{constrain()} function to squash all raw sensor readings into a uniform $0.0$ to $1.0$ decimal scale before applying the AI math.

\section{Detailed Firmware Code Explanation (\texttt{ESP32\_Air\_Quality.ino})}
The C++ firmware is engineered for industrial stability, running non-blocking loops and robust oversampling.

\subsection{Setup and Initialization}
The \texttt{setup()} function configures GPIO pins, runs an RGB LED hardware test, initializes the I2C OLED display, connects to Wi-Fi, and instantiates the local HTTP web server on port 80.

\subsection{The Continuous Loop (\texttt{loop()})}
The main loop utilizes a non-blocking 1.5-second timer (\texttt{millis()}) to ensure the microcontroller never freezes.

\subsubsection{Hardware Oversampling and Noise Reduction}
To mitigate electrical jitter, analog sensors are oversampled. The MQ gas sensors are read 24 times rapidly and averaged. The Dust sensor is read 8 times using specific microsecond delays to pulse its internal IR LED.

\subsubsection{Exponential Smoothing}
Even after oversampling, the firmware applies a commercial-grade exponential stabilization filter:
\begin{lstlisting}[language=C++]
livePM25 = (livePM25 * 0.85f) + (calculatedPM25 * 0.15f);
\end{lstlisting}
This mathematical filter ensures that new readings only shift the displayed value by 15\%, providing exceptionally smooth UI transitions on the web dashboard.

\subsubsection{Edge AI Execution and Actuator Control}
The firmware calls the Edge AI functions using the smoothed sensor data. Based on the returned Decision Tree Risk Class, it updates the physical RGB LED state. If the Logistic Regression hazard probability exceeds 85\%, it instantly triggers the buzzer alarm.

\subsubsection{Firebase Telemetry and Command Polling}
\begin{itemize}
    \item \textbf{Telemetry Push:} Every 3 seconds, the ESP32 constructs a JSON string of all variables and makes an HTTPS \texttt{PUT} request to the Firebase Realtime Database.
    \item \textbf{C2 Command Listener:} Every 1.5 seconds, the ESP32 executes an HTTPS \texttt{GET} request to Firebase to check for remote commands sent from the user's Web Dashboard (e.g., \texttt{buzzer\_on}, \texttt{oled\_off}). It verifies a timestamp to prevent executing duplicate commands.
\end{itemize}

\section{Conclusion}
The Smart Air Guardian's codebase represents a highly sophisticated amalgamation of embedded C++ logic, Python-based Machine Learning training, and modern web application development. By successfully shifting predictive inference to the edge while maintaining global cloud connectivity, the system offers an unprecedented level of real-time environmental awareness.

\end{document}
